\documentclass[11pt,letterpaper]{article}
\usepackage[margin=1in]{geometry}
\usepackage{amsmath,amssymb}
\usepackage{fancyhdr}
\usepackage{xcolor}
\usepackage{enumitem}
\usepackage[framemethod=tikz]{mdframed}
\usepackage[hidelinks]{hyperref}
\pagestyle{fancy}
\fancyhf{}
\lhead{\small AntsPhysicsLessons.com}
\rhead{\small Electricity \& Magnetism}
\cfoot{\small\thepage}
\renewcommand{\headrulewidth}{0.4pt}
\setlength{\parindent}{0pt}
\definecolor{keybg}{RGB}{240,244,250}
\newmdenv[backgroundcolor=keybg,linewidth=0pt,innerleftmargin=8pt,innerrightmargin=8pt,innertopmargin=6pt,innerbottommargin=6pt,skipabove=6pt]{answer}
\begin{document}
{\small\textsc{Unit 2 -- Gauss's Law}}\\[2pt]{\Large\bfseries 2.1 Fields of Continuous Charge Distributions}\\[8pt]Name: \rule{2.4in}{0.4pt} \hfill Date: \rule{1.4in}{0.4pt}\\[4pt]\hrule\vspace{10pt}{\small\itshape Placeholder worksheet for the site prototype. Free for classroom use.}\vspace{6pt}
\begin{enumerate}[leftmargin=*,itemsep=10pt]
\item A thin rod of length $L$ carries total charge $Q$ spread uniformly along it. Find the electric field at a point on the rod's axis a distance $a$ from the nearer end.
\vspace{2.2in}
\item A ring of radius $R$ carries charge $Q$ uniformly. Find the field on the axis a distance $z$ from the centre, and the value of $z$ where it is largest.
\vspace{2.2in}
\end{enumerate}
\end{document}